% Optional module: the source notation is that of (UZ4)--(UZ6).
% All statements below are exact; no localized ring is used.
\section{A short contact--conductor proof of the C1 memberships}

Write $d=\delta$ in this section; this local parameter is distinct from
the global Bell coefficient $d_{\rm B}$. Put
\[
 R_0=\mathbb Q[d,k,b,a,t],\qquad R=R_0[u,w,v],\qquad
 J=(u-k,w-dk,v-dk^2b^2).
\]
For $F\in R$, write $F_c=F(k,dk,dk^2b^2)$. Define
\[
 A=wv,\quad B=dv,\quad C=db^2uw,\quad D_0=d^2k^3b^2,
 \qquad M=(A,B,C),
\]
\[
 x=B-C,\qquad y=A-kC,\qquad z=uB-kC,\qquad h=kC-D_0.
\]
These are the original roots $g_0,h_{13},h_{23},h_{34}$; in particular,
$D_0=a_4^2$, whereas the source monomial in $g_4$ is $kD_0$.
Adjoin the separated radicals used in the original four-monomial circuit,
\[
 a_1=\sqrt{wv},\quad a_2=\sqrt{duv},\quad
 a_3=b\sqrt{dkuw},\quad a_4=dk^{3/2}b,
 \qquad q=a_1a_2-a_3a_4.
\]
Only $q^2\in R$ is needed in the scalar lift. The following three identities
are the entire conductor calculation:
\begin{align}
 C(u-k)&=z-ux,\label{CC-u}\\
 C(w-dk)&=dy-wx,\label{CC-w}\\
 C(v-dk^2b^2)&=b^2(wz+dky-kwx).\label{CC-v}
\end{align}
Thus $CJ\subset(x,y,z)$ without inverting any parameter.
For
\[
 r_*=y+z-h-2kx
\]
we also have the two exact identities
\begin{align}
 r_*&=(v-dk^2b^2)\bigl((w-dk)+d(u-k)\bigr)
       -dkb^2(u-k)(w-dk),\label{CC-star}\\
 kCr_*&=q^2-yz.\label{CC-repair}
\end{align}
The first says $r_*\in J^2$. The second supplies the otherwise missing
conductor factor through one of the original circuit squares.

\paragraph{Explicit monic division.}
Every $S\in R$ has a polynomial decomposition
\[
 S=S_c+(u-k)T_u+(w-dk)T_w+(v-dk^2b^2)T_v.                 \tag{CC1}
\]
One obtains it by substituting $u=k$, then $w=dk$, then $v=dk^2b^2$,
and using $X^n-Y^n=(X-Y)\sum_{j=0}^{n-1}X^{n-1-j}Y^j$ at each step.
There are no parameter denominators. If $S_c=0$, equations
\eqref{CC-u}--\eqref{CC-v} give
\[
 CS=x(-uT_u-wT_w-b^2kwT_v)
     +y(dT_w+dkb^2T_v)+z(T_u+b^2wT_v).                 \tag{CC2}
\]

\begin{lemma}[First-order and restricted second-order contact lifts]
In the above ring, the following assertions hold, with polynomial
coefficients in $R$:
\[
 M\cap J\subset (x,y,z),\qquad
 \bigl(M^2+(D_0^2)\bigr)\cap J^2
       \subset (x,y,z,h,q)^2.                         \tag{CC3}
\]
Here $(D_0^2)$ is the principal ideal generated by $D_0^2$; the second
assertion does not replace it by all mixed products involving $D_0$.
\end{lemma}
\begin{proof}
For the first assertion, write $F=R_AA+R_BB+R_CC$. Since
$A=y+kC$ and $B=x+C$,
\[
 F=R_Ay+R_Bx+CS,\qquad S=kR_A+R_B+R_C.
\]
The contact value is $F_c=C_cS_c$, where $C_c=d^2k^2b^2$ is a nonzero
polynomial. Since $R_0$ is a domain, $F_c=0$ implies $S_c=0$.
Apply (CC2). This cancellation proves a zero identity; it introduces no
coefficient $1/C_c$.

For the second assertion choose any displayed representation
$F=\sum_{i\le j}f_{ij}m_im_j+\rho D_0^2$, with
$(m_1,m_2,m_3)=(A,B,C)$. Substitute
\[
 A=y+kC,\qquad B=x+C,\qquad D_0=kC-h.
\]
This gives
\[
 F=Q_2+C(xH_x+yH_y+zH_z+hH_h)+C^2S,\qquad H_h=-2k\rho,       \tag{CC4}
\]
where $Q_2\in(x,y,h)^2$ and initially $H_z=0$.
The contact value again implies $S_c=0$. Use (CC2) on $CS$ and absorb
$C^2S$ into the four coefficients in (CC4). This leaves $H_h=-2k\rho$
unchanged and gives the same formula without $C^2S$.

The first derivatives of $F\in J^2$ vanish at contact. The first three
root gradients have determinant
\[
 \det\left(\frac{\partial(x,y,z)}{\partial(u,w,v)}\right)_c
       =-d^4k^4b^4\ne0.                              \tag{CC5}
\]
Equation~\eqref{CC-star} gives a gradient relation with coefficient vector
$(-2k,1,1,-1)$. Its coefficient at $h$ is $-1$.
Set $\lambda=2k\rho_c$. Subtract
$\lambda(-2k,1,1,-1)$ from $(H_x,H_y,H_z,H_h)$.
The last resulting contact value is zero, and the first three satisfy the
homogeneous linear system with determinant (CC5). Multiplication by its
adjugate and cancellation in the domain $R_0$ show that those three
contact values are zero too. No inverse of the determinant is used in a
coefficient: each resulting polynomial $J_i$ lies in $J$ and is divided
only by the monic procedure (CC1).
Consequently
\[
 F=Q_2+2k\rho_cCr_*+C\sum_{i=1}^4r_iJ_i
   =Q_2+2\rho_c(q^2-yz)+\sum_{i=1}^4r_i(CJ_i),          \tag{CC6}
\]
where $(r_1,r_2,r_3,r_4)=(x,y,z,h)$. Each $CJ_i$ is explicitly a
combination of $x,y,z$ by (CC2), which proves the assertion in the
unlocalized polynomial ring.
\end{proof}

\paragraph{Why the restriction matters.}
The unrestricted claim obtained by replacing $M^2+(D_0^2)$ with
$(A,B,C,D_0)^2$ is false. Indeed $Cr_*\in(A,B,C,D_0)^2\cap J^2$.
Under the positive weights below it has weight $42$. Every weight-$42$
member of the square of the original soft-root ideal must be divisible
by $x$, since $x$ has weight $20$ and every other original soft root has
weight at least $22$. But $Cr_*$ is not divisible by $x$: setting
$d=k=b=w=1$, $u=2$, $v=2$ gives $x=0$ and $Cr_*=2$.
The factor $k$ supplied by $D_0^2$ in (CC4) is therefore substantive.

\paragraph{Application to the complete source scalar and both normals.}
Use precisely the seven source coquotients and the definitions in
Appendix~\ref{app:c1-source-definitions}, renaming their clearing multiplier $\mathcal C=K^6Q^4T_n$
to distinguish it from the present conductor monomial $C$. In particular
\[
 K=1+k^2,\quad Q=K+ak^2b^2,\quad T_n=K^2-td^2k^4Qb^2,
 \quad D=1+u^2,\quad D_v=1+v^2,
\]
\[
 V_n=K^2-tQuvw,\quad Y_n=KT_nuw-avV_n,
 \quad B_0=vY_n-KT_n,
 \quad A_c=C_5T_n+aK(C_3D_v-C_5).
\]
The exact scalar $F_0$ has the following five-term form:
\begin{samepage}
\begin{align*}
F_0={}&4KQT_n^2[C_2(K^2-tQ)+C_bD_vtQ]D^2A^2\\
 &+4Q^2A_cD^2V_n^2B^2\\
 &+4K^3C_4DD_v(KT_nC-ab^2BV_n)^2\\
 &+4KQ^2C_aDB_0^2D_0^2\\
 &+2K^2QC_dD\{b^2B^2(Y_n^2-K^2T_n^2+2aKT_nV_n)
                         -2K^2T_n^2BC\}.
\end{align*}
\end{samepage}
Also the exact leading polynomial is
\[
 P=D(A^2+B^2)+K(C^2+D_0^2)-2K^2BC.
\]
Thus $E_0=F_0-8K^2D(K^6Q^4T_n)P$ belongs to
$M^2+(D_0^2)$ directly, without ideal division.

For the first normals put
\begin{align*}
 S_B={}&-8aK^2b^4C_4D_vV_n
          +8d^2Q^2k^6b^4C_aB_0+4KQb^2C_dB_0,\\
 S_C={}&8K^3T_nb^2C_4D_v.
\end{align*}
The normal numerator satisfies $S_n/d=BS_B+CS_C$. Hence
\begin{align*}
 L_1={}& A\,[8K^3QT_n^2(C_bD_v-C_2)D]\\
 &+B\,[-8dK^2Q^2uDvV_nA_c+dK^3auvS_B]
       +C\,[dK^3auvS_C],\\
 L_2={}&B\,[8K^2Q^2T_nDV_n(C_3D_v-C_5)-K^2T_nS_B]
       -C\,K^2T_nS_C.
\end{align*}
These formulas prove $L_1\in M$ and $L_2\in(B,C)$.
The native contact and derivative identities (SD9)--(SD11) in
Appendix~\ref{app:c1-source-definitions} prove
\[
 (E_0)_c=(\partial_uE_0)_c=(\partial_wE_0)_c
          =(\partial_vE_0)_c=0,\qquad (L_1)_c=(L_2)_c=0.      \tag{CC7}
\]
These are source polynomial identities, with no positivity hypothesis or
parameter inversion. Their proof uses only six short native balances and
the chain rule. The supplied small checker confirms the identities and
multiplies back the complete lifts from (CC2) and (CC6).
By (CC3),
\[
 E_0\in(x,y,z,h,q)^2,\qquad L_1,L_2\in(x,y,z).          \tag{CC8}
\]
Thus all three memberships previously expressed by the finite
Gr\"obner-basis witness have a single contact--conductor mechanism.

\paragraph{Bounded coefficients and the original Gram argument.}
Give $(d,k,b,a,t,u,w,v)$ weights $(2,2,6,0,0,2,4,18)$.
Then $A,y,z,h,q$ have weight $22$, while $B,C,x$ have weight $20$.
The three substitutions defining contact preserve weight, as do the
monic geometric-sum divisions after accounting for the weight of their
linear factor. The factored weight argument in
Appendix~\ref{app:c1-source-definitions} gives, without a large expansion,
\[
 \operatorname{wt}(E_0)\ge52,\quad
 \operatorname{wt}(L_1)\ge42,\quad
 \operatorname{wt}(L_2)\ge32.
\]
Consequently the construction gives lower bounds $8,20,10$ for the
three families of correction coefficients, respectively. The optional
exact checker obtains the stronger scalar minimum $56$ and coefficient
bounds $12,20,12$ for its particular monomial assignment.
All are positive. Every coefficient has nonnegative
exponents, including in the bounded shapes $a,t$; every correction
therefore has zero residue on the stated C1 branch. The full hard-normal
block and its positive residue are printed in
Appendix~\ref{app:c1-source-definitions}. Restoration is exactly that of the source theorem: restore
$\mathcal C=K^6Q^4T_n$ times the four positive packets subtracted there,
then undo the positive multiplier $\mathcal C\Delta^2\Lambda$. The present argument proves the complete
memberships on that branch and does not assert a global source cover.
